\documentclass[11pt,a4paper]{article}
\usepackage[T1]{fontenc}
\usepackage[utf8]{inputenc}
\usepackage{amsmath,amssymb,amsthm}
\usepackage[margin=1in]{geometry}
\usepackage[colorlinks=true,linkcolor=blue,urlcolor=blue]{hyperref}
\theoremstyle{plain}
\newtheorem{theorem}{Theorem}
\newtheorem{lemma}[theorem]{Lemma}
\newtheorem{proposition}[theorem]{Proposition}
\newtheorem{corollary}[theorem]{Corollary}
\theoremstyle{definition}
\newtheorem{remark}[theorem]{Remark}

\title{The empirical recurrence for OEIS A197360, proved by transfer matrix}
\author{Adrian Perez Fontelles\\ \small Independent researcher}
\date{1 September 2026}
\begin{document}
\maketitle

\begin{abstract}
OEIS A197360 counts $n\times4$ arrays over $\{0,\dots,4\}$ subject to a condition imposed on
every CELL over a neighbour set the entry names, and it carries an empirical recurrence of
order $53$ contributed by R.~H.~Hardin. It is true, and it is decidable rather than
empirical. A cell's neighbourhood meets three consecutive rows, so the state of a transfer
matrix is a pair of consecutive rows and each step settles the middle one; an array is a
walk, and $a(n)$ is a walk count on $S=2205$ vertices, which is $C$-finite. Whether the
conjectured recurrence annihilates it is then decided by a finite exact computation, with the
Cayley--Hamilton theorem bounding the work.
\end{abstract}

\noindent\small 2020 Mathematics Subject Classification. 05A15, 05B45, 15B36.\normalsize

\section{The conjecture}

OEIS A197360 is ``Number of nX4 0..4 arrays with each element x equal to the number its horizontal and vertical neighbors equal to 4,3,1,1,1 for x=0,1,2,3,4.''. It has offset $1$ and begins
\[
1, 17, 113, 803, 5855, 41459, 300943, 2224023,\ \dots
\]
The entry states, as a conjecture contributed by its author and never marked settled:
\begin{quote}
Empirical: a(n) = a(n-1) +25*a(n-2) +150*a(n-3) +414*a(n-4) -857*a(n-5) -9974*a(n-6) -42118*a(n-7) -67868*a(n-8) +291909*a(n-9) +1665950*a(n-10) +3428329*a(n-11) +810045*a(n-12) -23384598*a(n-13) -89246066*a(n-14) -146894491*a(n-15) +65409499*a(n-16) +914719497*a(n-17) +2219012323*a(n-18) +2086276830*a(n-19) -2516290285*a(n-20) -16891329383*a(n-21) -8932490107*a(n-22) +32645318584*a(n-23) -26461873134*a(n-24) -213227725865*a(n-25) -207243257310*a(n-26) -505793175949*a(n-27) +2898826155449*a(n-28) +7145836088511*a(n-29) -4629002872836*a(n-30) -2401405257780*a(n-31) -27307530798565*a(n-32) +59599301523601*a(n-33) -71872915971293*a(n-34) -134280224052192*a(n-35) +178887642807609*a(n-36) -470090079757467*a(n-37) +222856953175122*a(n-38) +1074246974800224*a(n-39) -596127798645363*a(n-40) +837976716887328*a(n-41) -652144155943248*a(n-42) -5992025098974012*a(n-43) +6264877490924760*a(n-44) +5326638416133696*a(n-45) -3126213675406656*a(n-46) -9240682701903840*a(n-47) +6180838932984192*a(n-48) +6548918601458304*a(n-49) -7829036052977664*a(n-50) +43487203235328*a(n-51) +1487460834902016*a(n-52) -303093721276416*a(n-53)
\end{quote}
As of the ``Last modified'' line on the live entry (Oct 03 2025 07:09:19, revision 9) this is still
recorded as empirical, and nothing on the entry records it as proved.

\section{Cells, not blocks}

The entry names the neighbours it means --- horizontal is left and right, vertical is up and
down, diagonal is northwest-to-southeast, antidiagonal is northeast-to-southwest, king-move is
all eight --- and every such set lies inside the $3\times3$ square centred on the cell. Here
the set is $\{(0,-1), (0,1), (-1,0), (1,0)\}$, written as offsets (row, column) from the cell. Every cell carries
the condition
\[
\text{a cell of value }x\text{ has exactly }x\text{ neighbours of the value paired with }x,
\]
and the count of neighbours enters it, so the boundary is not a detail that can be papered
over: a cell in the interior has $4$ neighbours in this set and a cell on the boundary
fewer, and a condition such as ``a strict majority of its neighbours'' therefore means
different things at the two. In particular one cannot pad the array with a border of zeros ---
that would give every cell the full $4$ and change what is being counted. The outside is
therefore carried
explicitly, as a sentinel row above the first and below the last and as a sentinel column
either side, and a neighbour that falls outside is simply absent from both counts.

Write an array as a sequence of rows $u_1,\dots,u_L$, each in
$\{0,\dots,4\}^{4}$. The neighbourhood of a cell in row $u_i$ meets $u_{i-1}$,
$u_i$ and $u_{i+1}$ and nothing else. Take as vertices the pairs
$(r,s)\in(\{0,\dots,4\}^{4}\cup\{\varnothing\})\times\{0,\dots,4\}^{4}$,
and put an edge $(r,s)\to(s,t)$ when row $s$ satisfies the condition at every one of its
$4$ cells, given $r$ above and $t$ below. An array with $L$ rows is then a walk of length
$L-1$ starting at $(\varnothing,u_1)$, each step settling one row, and the last row is
settled by the terminal weight, which evaluates it with $\varnothing$ below. Most of those pairs are indistinguishable, and exactly so. Every offset named above has $|dj|\le 1$, so whether the cell at column $j$ of the middle row is satisfied depends on the three rows only through their windows at $j-1,j,j+1$. Collect for each $j$ the indicator over the window of the row below, and call that tuple $P(r,s)$: the number of cells any $t$ violates is read off $P$, so is the number violated at the bottom edge, and the successor $(s,t)$ carries $P(s,t)$, which does not mention $r$. Hence two pairs $(r,s)$ and $(r',s)$ with the same $P$ have the same outgoing rows and the same successors, and the vertex is $(s,P)$. That leaves $S=2205$ vertices in place of $390625$. Hence,
with $M$ the adjacency matrix on $S=2205$ vertices, $\iota$ the starting vector and $\tau$ the
terminal weight,
\[
a(n)\;=\;\,\iota^{\!\top}M^{\,n-1}\tau ,
\]
the exponent fixed by the entry's own shape and checked against its published terms. In
particular $a$ is $C$-finite.

\section{The proof}

\begin{lemma}\label{lem:crit}
Let $q(t)=t^{53}-\sum_i c_it^{53-i}$ be the characteristic polynomial of the
conjectured recurrence. Then the residual $a(n)-\sum_ic_ia(n-i)$ equals
$\iota^{\!\top}M^{\,j}q(M)\tau$ for the corresponding walk index $j$, so the
recurrence holds from some point on if and only if $\iota^{\!\top}M^{j}q(M)\tau=0$ for all
large $j$; and it is enough to examine $j<S$.
\end{lemma}

\begin{proof}
Factor $M^{j}$ out of $M^{j+53}-\sum_ic_iM^{j+53-i}$. For the bound, $M^{S}$
is an integer combination of $I,M,\dots,M^{S-1}$ by Cayley--Hamilton, so the numbers
$u_j=\iota^{\!\top}M^{j}q(M)\tau$ satisfy the monic linear recurrence given by the
characteristic polynomial of $M$; $S$ consecutive zeros force every later one, and the last
nonzero $u_j$ pins the threshold exactly.
\end{proof}

\begin{theorem}
\[
a(n)\;=\;a(n-1) + 25\,a(n-2) + 150\,a(n-3) + 414\,a(n-4) - 857\,a(n-5) - 9974\,a(n-6) - 42118\,a(n-7) - 67868\,a(n-8) + 291909\,a(n-9) + 1665950\,a(n-10) + 3428329\,a(n-11) + 810045\,a(n-12) - 23384598\,a(n-13) - 89246066\,a(n-14) - 146894491\,a(n-15) + 65409499\,a(n-16) + 914719497\,a(n-17) + 2219012323\,a(n-18) + 2086276830\,a(n-19) - 2516290285\,a(n-20) - 16891329383\,a(n-21) - 8932490107\,a(n-22) + 32645318584\,a(n-23) - 26461873134\,a(n-24) - 213227725865\,a(n-25) - 207243257310\,a(n-26) - 505793175949\,a(n-27) + 2898826155449\,a(n-28) + 7145836088511\,a(n-29) - 4629002872836\,a(n-30) - 2401405257780\,a(n-31) - 27307530798565\,a(n-32) + 59599301523601\,a(n-33) - 71872915971293\,a(n-34) - 134280224052192\,a(n-35) + 178887642807609\,a(n-36) - 470090079757467\,a(n-37) + 222856953175122\,a(n-38) + 1074246974800224\,a(n-39) - 596127798645363\,a(n-40) + 837976716887328\,a(n-41) - 652144155943248\,a(n-42) - 5992025098974012\,a(n-43) + 6264877490924760\,a(n-44) + 5326638416133696\,a(n-45) - 3126213675406656\,a(n-46) - 9240682701903840\,a(n-47) + 6180838932984192\,a(n-48) + 6548918601458304\,a(n-49) - 7829036052977664\,a(n-50) + 43487203235328\,a(n-51) + 1487460834902016\,a(n-52) - 303093721276416\,a(n-53)
\]
for every $n>53$.
\end{theorem}

\begin{proof}
The vector $w=q(M)\tau$ was formed by $53$ matrix--vector products in exact integer
arithmetic and $\iota^{\!\top}M^{j}w$ evaluated for $j=0,1,\dots$, again exactly. Those
integers vanish from the index corresponding to $n=53+1$ onwards, and the run was
continued until the working vector was identically zero, which settles all larger $j$ at once.
By Lemma~\ref{lem:crit} the recurrence holds for every $n>53$.
\end{proof}

\section{Verification}

Three checks, each independent of the algebra above.

First, the transfer model reproduces all $20$ terms the entry publishes, exactly, in
integer arithmetic. That check earned its place here. The neighbour offsets are named in the
array's own frame --- horizontal means along a row --- but when the entry fixes the first
dimension, as in ``$2\times n$'', the walk runs along columns and the two components of every
offset swap. Sets such as king-move, or horizontal together with vertical, are unchanged by
that swap and hide the error completely; a set such as horizontal, diagonal and antidiagonal
is not, and the counts came out as those of a different member of the same family until the
offsets were transposed.

Second, the conjectured recurrence was evaluated directly on the entry's own published terms,
with no matrices involved, and holds wherever the proved range and the data overlap.

Third, the annihilation test of Lemma~\ref{lem:crit} was carried out in exact integer
arithmetic on the whole vector rather than on a sampled prefix.

\begin{thebibliography}{9}
\bibitem{oeis} The OEIS Foundation, \emph{The On-Line Encyclopedia of Integer
Sequences}, \url{https://oeis.org}, 2026. Sequence A197360.
\bibitem{stanley} R.~P.~Stanley, \emph{Enumerative Combinatorics}, Volume 1, 2nd ed.,
Cambridge University Press, 2011. (Transfer-matrix method, Section 4.7.)
\bibitem{fs} P.~Flajolet and R.~Sedgewick, \emph{Analytic Combinatorics}, Cambridge
University Press, 2009.
\bibitem{horn} R.~A.~Horn and C.~R.~Johnson, \emph{Matrix Analysis}, 2nd ed., Cambridge
University Press, 2013.
\end{thebibliography}

\end{document}
